\section{Carry Normalization}

Raw coefficients $c_k$ need not be valid base-$B$ digits. Carry normalization
defines
\[
  s_k = c_k + r_k,\qquad d_k = s_k \bmod B,\qquad r_{k+1} = \left\lfloor \frac{s_k}{B} \right\rfloor,
\]
with $r_0=0$. The normalized result is
\[
  XY = \sum_k d_k B^k.
\]

This step is nonlinear because thresholding and integer division redistribute
mass along the degree axis.

\subsection{Carry as Nonlinear Mass Flow}

After diagonal projection, the product is represented by a sequence of raw
coefficients
\[
    c_k = \sum_{i+j=k} a_i b_j.
\]
These coefficients need not be valid base-\(B\) digits. We therefore define
carry normalization by the recurrence
\[
    r_0 = 0,
\]
\[
    s_k = c_k + r_k,
\]
\[
    d_k = s_k \bmod B,
\]
\[
    r_{k+1} = \left\lfloor \frac{s_k}{B} \right\rfloor.
\]
Here \(r_k\) is the incoming carry at degree \(k\), \(d_k\) is the normalized
digit kept at degree \(k\), and \(r_{k+1}\) is the outgoing carry transported
from degree \(k\) to degree \(k+1\).

This interpretation views carry as a directed mass flow along the degree axis:
\[
    k \longrightarrow k+1.
\]
The mass transported across this edge is \(r_{k+1}\). Thus diagonal projection
is a bilinear operation, while carry normalization is a nonlinear flow process.

\begin{proposition}[Value preservation of carry normalization]
Let \(c=(c_k)\) be a finite sequence of nonnegative integer raw coefficients.
Let \(d=(d_k)\) be obtained from \(c\) by base-\(B\) carry normalization. Then
\[
    \sum_k c_k B^k = \sum_k d_k B^k.
\]
\end{proposition}

\begin{proof}
By the carry recurrence,
\[
    c_k + r_k = d_k + B r_{k+1}.
\]
Multiplying by \(B^k\), we obtain
\[
    c_k B^k + r_k B^k = d_k B^k + r_{k+1}B^{k+1}.
\]
Summing over \(k\), the carry terms telescope. Since \(r_0=0\) and the final
carry is eventually exhausted, we get
\[
    \sum_k c_k B^k = \sum_k d_k B^k.
\]
\end{proof}
